% GENERATED FILE. Edit canonical lessons, then run npm run generate:books.
% canonical-source-hash: fnv1a64:3ff064ed4dc10e39
% canonical-lessons: SA-C322-purvajah, SA-C322-vamsah, SA-C322-kulam, SA-C322-samsthapayati, SA-C322-pravartayati

\chapter{Ancestors}
\label{ch:sa-ancestors}
\hlchaptermodality{322}

\begin{chapteropening}
\textbf{By the end of this chapter:} \emph{I can tell the story of a family line.}

Complete the last lesson of chapter 322: I can tell the story of a family line.
\end{chapteropening}

\section[\texorpdfstring{\sk{पूर्वजः} (pūrvajaḥ)}{पूर्वजः (pūrvajaḥ)}]{\sk{पूर्वजः} (pūrvajaḥ) — an ancestor}

\label{lesson:SA-C322-purvajah}



\begin{quote}
Before the new one: say the Sanskrit for anoints, then the Sanskrit for pronounces.
\end{quote}

\subsection*{You'll want to know: \sk{पूर्वजः}}
\textbf{\sk{पूर्वजः}} — \emph{pūrvajaḥ} — \textquotedblleft{}an ancestor\textquotedblright{}.

\begin{quote}
\sk{मम} \sk{पूर्वजाः} \sk{तस्मिन्} \sk{ग्रामे} \sk{वसन्ति} \sk{स्म।}
\end{quote}

\begin{quote}
\emph{mama pūrvajāḥ tasmin grāme vasanti sma.}
\end{quote}

\begin{quote}
My ancestors used to live in that village.
\end{quote}

A masculine noun.

\begin{cousinweb}
\textbf{\sk{पूर्वजः}} (\emph{pūrvajaḥ}) is built from \emph{pūrva}, \textquotedblleft{}earlier, former\textquotedblright{}, and \emph{ja}, \textquotedblleft{}born\textquotedblright{}, from the root √jan; MW gives \emph{pūrvaja} as born before, and an ancestor or forefather (MW).
\end{cousinweb}

\begin{grammarlens}[title={What you've built}]
One more word for telling what happened.
\end{grammarlens}

\subsection*{Your turn}
\begin{itemize}
\raggedright
  \item \emph{Say it:} \emph{pūrvajaḥ}
  \item \emph{Say it:} \emph{pūrvajaḥ}, once more
  \item \emph{Say it:} \emph{mama pūrvajāḥ tasmin grāme vasanti sma.} — My ancestors used to live in that village.
  \item \emph{Recall:} say the Sanskrit for anoints, then the Sanskrit for pronounces, then say \emph{pūrvajaḥ} again
\end{itemize}

\begin{tcolorbox}[breakable,colback=teal!4,colframe=teal!35!black,title={Before you move on}]
What does \sk{पूर्वजः} mean? (\textquotedblleft{}An ancestor\textquotedblright{}.) Say it once more.
\end{tcolorbox}
\section[\texorpdfstring{\sk{वंशः} (vaṃśaḥ)}{वंशः (vaṃśaḥ)}]{\sk{वंशः} (vaṃśaḥ) — a family line; a dynasty}

\label{lesson:SA-C322-vamsah}



\begin{quote}
Before the new one: say the Sanskrit for pronounces, then the Sanskrit for an ancestor.
\end{quote}

\subsection*{You'll want to know: \sk{वंशः}}
\textbf{\sk{वंशः}} — \emph{vaṃśaḥ} — \textquotedblleft{}a family line; a dynasty\textquotedblright{}.

\begin{quote}
\sk{सः} \sk{राजा} \sk{प्राचीनस्य} \sk{वंशस्य} \sk{अस्ति।}
\end{quote}

\begin{quote}
\emph{saḥ rājā prācīnasya vaṃśasya asti.}
\end{quote}

\begin{quote}
That king comes from an ancient line.
\end{quote}

A masculine noun.

\begin{cousinweb}
\textbf{\sk{वंशः}} (\emph{vaṃśaḥ}) is built on the image of a bamboo cane: MW explains the sense of a lineage, race or dynasty from the line of joints in a bamboo (MW).
\end{cousinweb}

\begin{grammarlens}[title={What you've built}]
One more word for telling what happened.
\end{grammarlens}

\subsection*{Your turn}
\begin{itemize}
\raggedright
  \item \emph{Say it:} \emph{vaṃśaḥ}
  \item \emph{Say it:} \emph{vaṃśaḥ}, once more
  \item \emph{Say it:} \emph{saḥ rājā prācīnasya vaṃśasya asti.} — That king comes from an ancient line.
  \item \emph{Recall:} say the Sanskrit for pronounces, then the Sanskrit for an ancestor, then say \emph{vaṃśaḥ} again
\end{itemize}

\begin{tcolorbox}[breakable,colback=teal!4,colframe=teal!35!black,title={Before you move on}]
What does \sk{वंशः} mean? (\textquotedblleft{}A family line; a dynasty\textquotedblright{}.) Say it once more.
\end{tcolorbox}
\section[\texorpdfstring{\sk{कुलम्} (kulam)}{कुलम् (kulam)}]{\sk{कुलम्} (kulam) — a clan; a family}

\label{lesson:SA-C322-kulam}



\begin{quote}
Before the new one: say the Sanskrit for an ancestor, then the Sanskrit for a family line.
\end{quote}

\subsection*{You'll want to know: \sk{कुलम्}}
\textbf{\sk{कुलम्}} — \emph{kulam} — \textquotedblleft{}a clan; a family\textquotedblright{}.

\begin{quote}
\sk{तस्य} \sk{कुलं} \sk{प्रसिद्धम्} \sk{अस्ति।}
\end{quote}

\begin{quote}
\emph{tasya kulaṃ prasiddham asti.}
\end{quote}

\begin{quote}
His clan is famous.
\end{quote}

A neuter noun.

\begin{cousinweb}
\textbf{\sk{कुलम्}} (\emph{kulam}) MW gives \emph{kula} as a herd, troop or multitude, and also a race, family, community or noble house, without a derivation, so its origin is uncertain (MW).
\end{cousinweb}

\begin{grammarlens}[title={What you've built}]
One more word for telling what happened.
\end{grammarlens}

\subsection*{Your turn}
\begin{itemize}
\raggedright
  \item \emph{Say it:} \emph{kulam}
  \item \emph{Say it:} \emph{kulam}, once more
  \item \emph{Say it:} \emph{tasya kulaṃ prasiddham asti.} — His clan is famous.
  \item \emph{Recall:} say the Sanskrit for an ancestor, then the Sanskrit for a family line, then say \emph{kulam} again
\end{itemize}

\begin{tcolorbox}[breakable,colback=teal!4,colframe=teal!35!black,title={Before you move on}]
What does \sk{कुलम्} mean? (\textquotedblleft{}A clan; a family\textquotedblright{}.) Say it once more.
\end{tcolorbox}
\section[\texorpdfstring{\sk{संस्थापयति} (saṃsthāpayati)}{संस्थापयति (saṃsthāpayati)}]{\sk{संस्थापयति} (saṃsthāpayati) — founds; establishes}

\label{lesson:SA-C322-samsthapayati}



\begin{quote}
Before the new one: say the Sanskrit for a family line, then the Sanskrit for a clan.
\end{quote}

\subsection*{You'll want to know: \sk{संस्थापयति}}
\textbf{\sk{संस्थापयति}} — \emph{saṃsthāpayati} — \textquotedblleft{}founds; establishes\textquotedblright{}.

\begin{quote}
\sk{राजा} \sk{नवं} \sk{नगरं} \sk{संस्थापयति।}
\end{quote}

\begin{quote}
\emph{rājā navaṃ nagaraṃ saṃsthāpayati.}
\end{quote}

\begin{quote}
The king founds a new city.
\end{quote}

Past (he …): \sk{सः} \sk{नगरं} \sk{संस्थापितवान्}, \emph{saḥ nagaraṃ saṃsthāpitavān}.

\begin{cousinweb}
\textbf{\sk{संस्थापयति}} (\emph{saṃsthāpayati}) is built from \emph{sam} and the causative of the root √sthā, \textquotedblleft{}to stand\textquotedblright{}, as in \emph{sthāpayati}; MW gives \emph{saṃ-sthāpayati} as to make stand firm, raise up, restore (MW).
\end{cousinweb}

\begin{grammarlens}[title={What you've built}]
One more word for telling what happened.
\end{grammarlens}

\subsection*{Your turn}
\begin{itemize}
\raggedright
  \item \emph{Say it:} \emph{saṃsthāpayati}
  \item \emph{Say it:} \emph{saṃsthāpayati}, once more
  \item \emph{Say it:} \emph{rājā navaṃ nagaraṃ saṃsthāpayati.} — The king founds a new city.
  \item \emph{Recall:} say the Sanskrit for a family line, then the Sanskrit for a clan, then say \emph{saṃsthāpayati} again
\end{itemize}

\begin{tcolorbox}[breakable,colback=teal!4,colframe=teal!35!black,title={Before you move on}]
What does \sk{संस्थापयति} mean? (\textquotedblleft{}Founds; establishes\textquotedblright{}.) Say it once more.
\end{tcolorbox}
\section[\texorpdfstring{\sk{प्रवर्तयति} (pravartayati)}{प्रवर्तयति (pravartayati)}]{\sk{प्रवर्तयति} (pravartayati) — starts; introduces}

\label{lesson:SA-C322-pravartayati}



\begin{quote}
Before the new one: say the Sanskrit for a clan, then the Sanskrit for founds.
\end{quote}

\subsection*{You'll want to know: \sk{प्रवर्तयति}}
\textbf{\sk{प्रवर्तयति}} — \emph{pravartayati} — \textquotedblleft{}starts; introduces\textquotedblright{}.

\begin{quote}
\sk{राजा} \sk{नवम्} \sk{उत्सवं} \sk{प्रवर्तयति।}
\end{quote}

\begin{quote}
\emph{rājā navam utsavaṃ pravartayati.}
\end{quote}

\begin{quote}
The king starts a new festival.
\end{quote}

Past (he …): \sk{सः} \sk{उत्सवं} \sk{प्रवर्तितवान्}, \emph{saḥ utsavaṃ pravartitavān}.

\begin{cousinweb}
\textbf{\sk{प्रवर्तयति}} (\emph{pravartayati}) is built from \emph{pra} and the causative of the root √vṛt, \textquotedblleft{}to turn, roll\textquotedblright{}; MW gives \emph{pravartita} as set in motion, set up, established, introduced (MW).
\end{cousinweb}

\begin{grammarlens}[title={What you've built}]
One more word for telling what happened.
\end{grammarlens}

\subsection*{Your turn}
\begin{itemize}
\raggedright
  \item \emph{Say it:} \emph{pravartayati}
  \item \emph{Say it:} \emph{pravartayati}, once more
  \item \emph{Say it:} \emph{rājā navam utsavaṃ pravartayati.} — The king starts a new festival.
  \item \emph{Recall:} say the Sanskrit for a clan, then the Sanskrit for founds, then say \emph{pravartayati} again
\end{itemize}

\begin{tcolorbox}[breakable,colback=teal!4,colframe=teal!35!black,title={Before you move on}]
What does \sk{प्रवर्तयति} mean? (\textquotedblleft{}Starts; introduces\textquotedblright{}.) Say it once more.
\end{tcolorbox}
